% sphere-heat-transfer.tex

\newpage
\section{Heat transfer to a sphere in equilibrium air}
\label{sphere-heat-transfer-sec}
\index{gas model!look-up table!example of use}
%
This example continues the modelling of hypersonic flow over blunt bodies 
and looks at the heat transfer to a spherical probe \cite{rose_stark_58} in high temperature equilibrium air. 
It takes use of the Python language further by automating the process of running 
a simulation, adjusting the grid and then running a subsequent simulation on the
the adjusted grid.
The specific input file for each stage of the overall simulation is constructed from a template
in which a few parameters are left unspecified.
Most of the effort has gone into the coordinating script which has functions
for running stages of the simulation as subprocesses and also has functions
which fit a Bezier curve to the shock located in the flow field.

\begin{figure}[htbp]
\begin{center}
\includegraphics[width=10cm,viewport=23 70 296 351,clip=true]{../2D/sphere-heat-transfer/sphere0-layout.pdf}
\end{center}
\caption{Schematic diagram of the geometry for a sphere wrapped by a SuperBlock2D grid.}
\label{sphere-heat-transfer-geometry-fig}
\end{figure}

\medskip
The original experiments used a probe with a spherical nose, located in a small shock tube.
The free-stream flow was initiated with the arrival of a strong shock and the useful test period 
in the experiments was terminated with the arrival of driver gas.
From Figure\,12 in Rose and Detra's paper \cite{rose_stark_58}, we choose the point corresponding to
$p_1 = 1$\,cm Hg (1.33\,kPa) and $M_s$=8 which has a stagnation-point heat transfer of $30\pm2.0\,{\rm MW/m}^2$.
To keep the grid resolution requirements small,
we will start with an initial test gas pressure $p_1 = 6.7$\,Pa much lower 
than that used in the original experiments.  
Assuming that the chemistry doesn't change too much with the change in pressure, we can scale the stagnation-point
heat transfer as $\dot{q}_{s-sim} = \left( \frac{p_{1-sim}}{p_{1-expt}} \right)^{0.5} \dot{q}_{s-expt}$ 
to get an expected value of $2.212\pm0.14\,{\rm MW/m}^2$ for our low-pressure simulation.

\medskip
For a $M_s = 8$ incident shock in air at 296\,K, the post-shock, free-stream conditions
are $p_{\infty} = 535.6$\,kPa, $T_{\infty} = 2573.5$\,K, and $u_{\infty} = 2436.5$\,m/s.
This assumes fully-equilibrium chemistry for the gas.
The snap-shots of results for the staged simulation are shown in Figures \ref{sphere-heat-transfer-stage-0-fig} through
\ref{sphere-heat-transfer-stage-4-fig} which show the temperature field at the end of each stage
and the mesh used for that stage. 

\begin{figure}[htbp]
\begin{center}
\mbox{
\includegraphics[width=0.4\textwidth]{../2D/sphere-heat-transfer/sphere0-T-field.png}
\includegraphics[width=0.4\textwidth]{../2D/sphere-heat-transfer/sphere0-mesh.png}
}
\end{center}
\caption{Temperature field and mesh for stage 0.  
  The control points for the Bezier curve have been set so as to accommodate 
  a shock in ideal (nonreacting) air and the clustering is fairly strong so that
  the boundary layer on the sphere surface may be resolved.
  The wall-clock time required to run this simulation 10 body lengths (27$\mu$s) is 23 seconds 
  on 4 processors of \texttt{geyser}.
  10103 time steps were made and the size of the time step at the end of the simulation is 2.479\,ns.
  At the end of the simulation, the estimated value of stagnation-point heat transfer is
  $\dot{q}_s=2.156\,{\rm MW/m}^2$ and the cell Reynolds number at the stagnation point is
  ${\rm Re}_{wall} = \frac{\rho_{wall} a_{wall} \Delta x}{\mu_{wall}} = 3.85$.
  Here $\Delta x$ is the width of the cell out from the wall.}
\label{sphere-heat-transfer-stage-0-fig}
\end{figure}

\begin{figure}[htbp]
\begin{center}
\mbox{
\includegraphics[width=0.4\textwidth]{../2D/sphere-heat-transfer/sphere1-T-field.png}
\includegraphics[width=0.4\textwidth]{../2D/sphere-heat-transfer/sphere1-mesh.png}
}
\end{center}
\caption{Temperature field and mesh for stage 1.
  The Bezier points have been adapted to the shock from stage 0 but the number of cells
  in each direction remains at 20$\times$20, as for stage 0.
  With the finer cells, the size of the time step decreased and this stage required 55 seconds of wall-clock time
  to extend the simulation a further 10 body lengths in 23125 time steps.
  $\dot{q}_s=2.260\,{\rm MW/m}^2$ and ${\rm Re}_{wall} =2.94$}
\label{sphere-heat-transfer-stage-1-fig}
\end{figure}

\begin{figure}[htbp]
\begin{center}
\mbox{
\includegraphics[width=0.4\textwidth]{../2D/sphere-heat-transfer/sphere2-T-field.png}
\includegraphics[width=0.4\textwidth]{../2D/sphere-heat-transfer/sphere2-mesh.png}
}
\end{center}
\caption{Temperature field and mesh for stage 2.
  The Bezier points have not been adapted further for this stage but the number of cells
  has been increased to 30$\times$30.
  The size of the time step decreased further and this stage required 130 seconds of wall-clock time
  to extend the simulation only 5 body lengths (13.5\,$\mu$s) in 24293 time steps.
  $\dot{q}_s=2.257\,{\rm MW/m}^2$ and ${\rm Re}_{wall} =2.39$}
\label{sphere-heat-transfer-stage-2-fig}
\end{figure}

\begin{figure}[htbp]
\begin{center}
\mbox{
\includegraphics[width=0.4\textwidth]{../2D/sphere-heat-transfer/sphere3-T-field.png}
\includegraphics[width=0.4\textwidth]{../2D/sphere-heat-transfer/sphere3-mesh.png}
}
\end{center}
\caption{Temperature field and mesh for stage 3.
  The Bezier points have been adapted to the shock from stage 2 
  and the cells have been increased to 40$\times$40.
  The size of the time step is now 0.319\,ns and this stage required 469 seconds of wall-clock time
  to extend the simulation only 5 body lengths (13.5\,$\mu$s) in 42660 time steps.
  $\dot{q}_s = 2.260\,{\rm MW/m}^2$ and ${\rm Re}_{wall} =1.99$}.
\label{sphere-heat-transfer-stage-3-fig}
\end{figure}

\begin{figure}[htbp]
\begin{center}
\mbox{
\includegraphics[width=0.4\textwidth]{../2D/sphere-heat-transfer/sphere4-T-field.png}
\includegraphics[width=0.4\textwidth]{../2D/sphere-heat-transfer/sphere4-mesh.png}
}
\end{center}
\caption{Temperature field and mesh for stage 4.
  The Bezier points have not been further adapted but the number of cells has been
  increased to 80$\times$80 to test the sensitivity of the heat transfer estimate.
  The size of the time step is now 0.086\,ns and this stage required 8950 seconds of wall-clock time
  to extend the simulation a further 5 body lengths (13.5\,$\mu$s) in 157420 time steps.
  $\dot{q}_s =2 .217\,{\rm MW/m}^2$ and ${\rm Re}_{wall} = 1.25$}.
\label{sphere-heat-transfer-stage-4-fig}
\end{figure}

Figure \ref{sphere-norm-heat-transfer-fig} shows the distribution of heat transfer around the nose
compared with the experimental data reported Kemp, Rose and Detra \cite{kemp_etal_1959}.
In the simulation data, there are small disturbances at the corners of blocks 
(at approximately 20 degrees and then again approaching 90 degrees) but they are quite small.

\begin{figure}[htbp]
\begin{center}
\includegraphics[width=12cm]{../2D/sphere-heat-transfer/sphere_norm_heat_transfer.pdf}
\end{center}
\caption{Temperature around the sphere for stage 4.  
   The experimental data is from Kemp, Rose and Detra \cite{kemp_etal_1959}.}.
\label{sphere-norm-heat-transfer-fig}
\end{figure}

\clearpage

\subsection{Template input script (.py)}\index{ExistingSolution!example of use}\index{univariate function!RobertsClusterFunction!example of use}
\topbar
\lstinputlisting[language={}]{../2D/sphere-heat-transfer/sphere.input.template}
\bottombar

\subsection{Coordinating script (.py)}
\topbar
\lstinputlisting[language={}]{../2D/sphere-heat-transfer/run_adaptive_simulation.py}
\bottombar

\newpage
\subsection{Shell script for postprocessing}
\topbar
\lstinputlisting[language={}]{../2D/sphere-heat-transfer/plot_heat_transfer.sh}
\bottombar

\newpage
\subsection{Notes}
\begin{itemize}
\item The look-up table for the equilibrium air equation of state is set up 
   as for the Sawada sphere example.
\end{itemize}
